\begin{problem}{\bf Meet-in-the-middle attack on Double-DES} (10 points)\\
In this question, you have to assume that the total amount of storage
space available is 1,000 Terabytes ($\approx 2^{53}$ bits).
\begin{enumerate}[label=(\alph*)]
\item Describe a known-plaintext meet-in-the-middle attack on Double-DES
that recovers the secret key. You should justify the correctness
of the algorithm, estimate the expected number of DES operations
the algorithm takes, and justify that the algorithm requires at most
1,000 Terabytes of storage. (A \emph{DES operation} is a
single-DES encryption or a single-DES decryption.) Of course,
you should design your algorithm to be as fast as you can.

\item Suppose that you have access to a network of 1,000 computers, each
computer capable of performing 5,000,000 DES operations per second.
The computers have access to shared storage, and you can assume
for simplicity that accessing this storage takes negligible time.
Estimate the expected number of calendar years your algorithm takes
when run on this network. Compare this to the estimated number of
calendar years that exhaustive key search on Double-DES would take
when run on the same network.
\end{enumerate}
\end{problem}


\begin{loigiai}
\textbf{(a) Tấn công Meet-in-the-Middle trên Double-DES}

Gọi khóa bí mật Double-DES là $(k_1,k_2)$, trong đó
$k_1,k_2 \in \{0,1\}^{56}$. Quá trình mã hóa được xác định bởi
\[
c = E_{k_2}(E_{k_1}(p)),
\]
trong đó $E_k$ và $D_k=E_k^{-1}$ lần lượt là phép mã hóa
và giải mã DES với khóa $k$.

Giả sử ta biết ba cặp plaintext/ciphertext
$(p_1,c_1)$, $(p_2,c_2)$ và $(p_3,c_3)$.

\textbf{1. Mô tả thuật toán}

Để đáp ứng giới hạn bộ nhớ, ta chia không gian khóa $h_2$
thành $2^{10}$ nhóm, mỗi nhóm gồm $2^{46}$ khóa. Với mỗi
nhóm, thực hiện thuật toán sau.

\begin{algorithm}[H]
\caption{Tấn công Meet-in-the-Middle trên Double-DES}
\KwIn{Ba cặp $(p_1,c_1)$, $(p_2,c_2)$, $(p_3,c_3)$}
\KwOut{Khóa bí mật $(k_1,k_2)$}

Chia $\{0,1\}^{56}$ thành $2^{10}$ nhóm $G$, mỗi nhóm có
$2^{46}$ khóa\;

\ForEach{nhóm $G$}{
    Khởi tạo bảng $T$\;

    \ForEach{$h_2 \in G$ \tcp*{$h_2$ là một giá trị dự đoán của $k_2$}}{
        $x \gets D_{h_2}(c_1)$\;
        Lưu $(x,\operatorname{index}(h_2))$ vào bảng $T$\;
    }

    Sắp xếp bảng $T$ theo thành phần $x$\;

    \ForEach{$h_1 \in \{0,1\}^{56}$\tcp*{$h_1$ là một giá trị dự đoán của $k_1$}}{
        $y \gets E_{h_1}(p_1)$\;
        Tìm tất cả mục trong $T$ có thành phần thứ nhất bằng $y$\;

        \ForEach{mục khớp tương ứng với khóa $h_2$}{
            \If{$E_{h_2}(E_{h_1}(p_2))=c_2$}{
                \If{$E_{h_2}(E_{h_1}(p_3))=c_3$}{
                    Xuất $(h_1,h_2)$ và dừng\;
                }
            }
        }
    }
}
\end{algorithm}

Ở đây, $\operatorname{index}(h_2)$ là chỉ số 46 bit xác định
khóa $h_2$ trong nhóm hiện tại. Sau khi tìm được chỉ số,
ta khôi phục khóa $h_2$ tương ứng.

\textbf{2. Chứng minh tính đúng đắn}

Từ công thức mã hóa Double-DES, ta có
\[
c_1=E_{k_2}(E_{k_1}(p_1)).
\]
Áp dụng phép giải mã $D_{k_2}$ lên hai vế, suy ra
\[
D_{k_2}(c_1)=E_{k_1}(p_1).
\]

Do đó, khi xử lý nhóm chứa khóa bí mật $k_2$, giá trị
$D_{k_2}(c_1)$ trong bảng sẽ trùng với giá trị
$E_{k_1}(p_1)$ khi thuật toán xét $h_1=k_1$. Vì vậy,
khóa bí mật thực sự luôn được đưa vào tập ứng viên.

Thuật toán tiếp tục kiểm tra hai cặp plaintext/ciphertext
còn lại:
\[
E_{h_2}(E_{h_1}(p_2))=c_2,
\]
\[
E_{h_2}(E_{h_1}(p_3))=c_3.
\]

Khóa bí mật thực sự luôn thỏa mãn cả hai điều kiện trên.
Theo giả định các kết quả mã hóa có tính chất ngẫu nhiên,
với một khóa sai, xác suất khớp cả ba cặp đã biết xấp xỉ
$2^{-192}$. Vì có khoảng $2^{112}$ cặp khóa, số khóa sai
kỳ vọng vượt qua cả ba phép kiểm tra là
\[
F_K \approx \frac{2^{112}}{2^{192}}=2^{-80}.
\]
Do đó, xác suất thu được khóa sai là rất nhỏ.

\textbf{3. Ước tính số phép toán DES}

\begin{itemize}
    \item Có $2^{10}$ nhóm khóa $h_2$, mỗi nhóm chứa $2^{46}$ khóa.
    Với mỗi nhóm, việc xây dựng bảng cần $2^{46}$ phép giải mã DES,
    còn việc duyệt toàn bộ $2^{56}$ khóa $h_1$ cần $2^{56}$ phép mã hóa DES.
    Do đó, tổng chi phí của hai công việc này là
    \[
    T_{\mathrm{main}}
    =2^{10}\cdot 2^{46}+2^{10}\cdot 2^{56}
    =2^{56}+2^{66}.
    \]

    \item Theo giả định các giá trị trung gian phân bố ngẫu nhiên,
    số cặp khóa $(h_1,h_2)$ thỏa mãn
    \[
    E_{h_2}(E_{h_1}(m_1))=c_1
    \]
    có kỳ vọng xấp xỉ
    \[
    \mathbb{E}[X_1]
    \approx \frac{2^{112}}{2^{64}}
    =2^{48}.
    \]
    Mỗi ứng viên cần $2$ phép toán DES để kiểm tra cặp plaintext/ciphertext
    thứ hai. Vì vậy, kỳ vọng của chi phí kiểm tra là
    \[
    T_{\mathrm{check},2}=2\cdot 2^{48}=2^{49}.
    \]

    \item Theo công thức số khóa sai kỳ vọng, sau khi kiểm tra hai cặp
    bản rõ--bản mã, số khóa sai còn lại có kỳ vọng
    \[
    \mathbb{E}[F_K]
    \approx \frac{2^{112}}{2^{64\cdot 2}}
    =2^{-16}.
    \]
    Cộng với khóa đúng, số ứng viên cần kiểm tra bằng cặp thứ ba
    có kỳ vọng xấp xỉ $1+2^{-16}$. Do mỗi ứng viên cần $2$ phép toán
    DES, kỳ vọng của chi phí kiểm tra là
    \[
    T_{\mathrm{check},3}=2(1+2^{-16}).
    \]
\end{itemize}

Vậy tổng số phép toán DES kỳ vọng là
\[
\begin{aligned}
\mathbb{E}[T]
&=T_{\mathrm{main}}+T_{\mathrm{check},2}
  +T_{\mathrm{check},3}\\
&=2^{56}+2^{66}+2^{49}+2(1+2^{-16})\\
&\approx \boxed{2^{66}\text{ phép toán DES}}.
\end{aligned}
\]

Các chi phí kiểm tra ứng viên đã được tính vào tổng số phép toán.
Tuy nhiên, chúng nhỏ hơn nhiều so với chi phí duyệt khóa $2^{66}$,
nên không làm thay đổi bậc độ lớn của tổng chi phí kỳ vọng.

\textbf{4. Phân tích dung lượng lưu trữ}

Bảng $T$ chứa $2^{46}$ mục. Mỗi mục lưu một giá trị trung
gian 64 bit và một chỉ số khóa 46 bit, tổng cộng 110 bit.

Vì vậy, dung lượng lưu trữ cần thiết là
\[
S=110\cdot 2^{46}\text{ bit}.
\]

Đổi trực tiếp sang TB theo hệ đơn vị thập phân,
với $1\text{ TB}=10^{12}$ byte:
\[
\begin{aligned}
S
&=\frac{110\cdot 2^{46}}{8\cdot 10^{12}}\text{ TB}\\
&\approx \boxed{967.6\text{ TB}<1000\text{ TB}}.
\end{aligned}
\]

Vì thuật toán chỉ lưu bảng của một nhóm tại mỗi thời điểm,
nó đáp ứng giới hạn bộ nhớ đã cho trong mô hình lưu trữ
lý thuyết, chưa tính bộ nhớ phụ trợ của cấu trúc dữ liệu.

\textbf{(b) Ước tính thời gian chạy trên mạng máy tính}

Mạng có $1{,}000$ máy tính, mỗi máy thực hiện được
$5{,}000{,}000$ phép toán DES mỗi giây. Tổng tốc độ xử lý là
\[
R=1000\cdot5\times10^6
=5\times10^9\text{ phép toán DES/giây}.
\]

Một năm có khoảng
\[
D=365.25\cdot24\cdot3600
=31{,}557{,}600\text{ giây}.
\]

Theo câu (a), kỳ vọng số phép toán DES của thuật toán
Meet-in-the-Middle là
\[
\begin{aligned}
\mathbb{E}[T]
&=2^{56}+2^{66}+2^{49}+2(1+2^{-16}).
\end{aligned}
\]

Do đó, thời gian chạy trung bình là
\[
\begin{aligned}
Y_{\mathrm{MITM}}
&=\frac{\mathbb{E}[T]}{R\cdot D}\\
&=\frac{2^{56}+2^{66}+2^{49}+2(1+2^{-16})}
{5\times10^9\cdot31{,}557{,}600}\\
&\approx\boxed{468\text{ năm}}.
\end{aligned}
\]

Đối với vét cạn khóa Double-DES, có $2^{112}$ cặp khóa.
Trong trường hợp khóa đúng được phân bố đều theo thứ tự thử,
khi đó kỳ vọng số cặp khóa cần thử là $2^{111}$. Mỗi cặp khóa
cần tối đa $2$ phép toán DES để kiểm tra, nên kỳ vọng số phép toán
DES xấp xỉ $2^{112}$.

Vì vậy,  thời gian vét cạn trung bình là
\[
\begin{aligned}
Y_{\mathrm{brute-force}}
&=\frac{2^{112}}{R\cdot D}\\
&\approx\boxed{3.29\times10^{16}\text{ năm}}.
\end{aligned}
\]

Như vậy, thuật toán Meet-in-the-Middle cần khoảng trung bình $468$ năm,
nhanh hơn rất nhiều so với vét cạn khóa Double-DES, vốn cần
khoảng trung bình $3.29\times10^{16}$ năm.

\end{loigiai}
